\documentclass[10pt]{article}
\usepackage[T1]{fontenc}
\usepackage[utf8]{inputenc}
\usepackage[french]{babel}
\usepackage{amsmath,amssymb,mathrsfs,mathpazo}
\usepackage{geometry}
\geometry{a4paper,margin=2cm}
\usepackage{graphicx,array,tabularx,tabularray,colortbl,multirow,multicol,booktabs,enumitem,pifont,ulem,fancybox,diagbox,cancel}
\usepackage{pstricks,pstricks-add,pst-plot,pst-tree,pst-node,pst-text,pst-func,pst-math,pst-bspline,pst-3dplot,pst-eucl,pst-solides3d}
\let\pscalculate\undefined
\usepackage{pst-bezier}
\usepackage{tikz,pgfplots,tkz-tab}
\usetikzlibrary{arrows,arrows.meta,patterns,calc,positioning,trees,decorations.markings,shapes.geometric,shapes.misc,shapes.arrows,intersections,decorations.pathmorphing,decorations.pathreplacing,backgrounds,shadows,babel}
\pgfplotsset{compat=1.18}
\usepgfplotslibrary{fillbetween}
\usepackage[np]{numprint}
\usepackage{listings,makecell,pas-tableur,twoopt}
\let\barre\undefined
\usepackage[pstricks]{tablvar}
\definecolor{darkgreen}{RGB}{0,100,0}
\providecommand{\red}{\color{red}}
\providecommand{\blue}{\color{blue}}
\providecommand{\green}{\color{green}}
\providecommand{\black}{\color{black}}
\providecommand{\white}{\color{white}}
\providecommand{\gray}{\color{gray}}
\providecommand{\cyan}{\color{cyan}}
\providecommand{\magenta}{\color{magenta}}
\providecommand{\yellow}{\color{yellow}}
\providecommand{\violet}{\color{violet}}
\providecommand{\orange}{\color{orange}}
\providecommand{\euro}{\texteuro}
\providecommand{\EUR}{\texteuro}
\providecommand{\ieme}{\textsuperscript{e}}
\providecommand{\degres}{\ensuremath{^\circ}}
\providecommand{\degree}{\ensuremath{^\circ}}
\providecommand{\vv}[1]{\overrightarrow{#1}}
\providecommand{\Coord}[1]{#1}
\providecommand{\decosix}{}
\providecommand{\starredbullet}{\ensuremath{\bullet}}
\pagestyle{empty}
\setlength{\parindent}{0pt}
\usepackage{adjustbox}\DeclareUnicodeCharacter{2236}{:}\DeclareUnicodeCharacter{2190}{\ensuremath{\leftarrow}}
\providecommand{\R}{\mathbb{R}}
\providecommand{\N}{\mathbb{N}}
\providecommand{\Z}{\mathbb{Z}}
\providecommand{\Q}{\mathbb{Q}}
\providecommand{\D}{\mathbb{D}}
\providecommand{\C}{\mathbb{C}}
\providecommand{\e}{{\mathrm{e}}}
\providecommand{\vect}[1]{\overrightarrow{#1}}
\providecommand{\vectt}[1]{\overrightarrow{\mathrm{#1}}}
\providecommand{\vvt}[1]{\overrightarrow{\mathrm{#1}}}
\providecommand{\np}[2][]{{#2~#1}}
\providecommand{\vv}[1]{{\overrightarrow{#1}}}
\renewcommand{\d}{{\mathrm{d}}}
\providecommand{\degres}{{^\circ}}
\providecommand{\degre}{{^\circ}}
\providecommand{\blue}{\color{blue}}
\providecommand{\red}{\color{red}}
\providecommand{\enskip}{\quad}
\providecommand{\llbracket}{\lbrack\!\lbrack}
\providecommand{\rrbracket}{\rbrack\!\rbrack}
\providecommand{\euro}{{\text{€}}}
\providecommand{\no}{{\text{n°}}}
\providecommand{\interval}[2]{{[#1;#2]}}
\providecommand{\intervFO}[2]{{[#1;#2[}}
\providecommand{\dotfill}{\ldots}

\providecommand{\ds}{\displaystyle}
\providecommand{\textding}[1]{#1}
\providecommand{\Oij}{\(\mathrm{O};\vec{\imath},\vec{\jmath}\)}
\providecommand{\Oijk}{\(\mathrm{O};\vec{\imath},\vec{\jmath},\vec{k}\)}
\providecommand{\Ouv}{\(\mathrm{O};\vec{u},\vec{v}\)}
\providecommand{\intervFF}[2]{\ensuremath{[#1;#2]}}
\renewcommand{\euro}{\ifmmode\text{\texteuro}\else\texteuro\fi}
\DeclareUnicodeCharacter{00B0}{\ifmmode^{\circ}\else\textdegree\fi}
\DeclareUnicodeCharacter{20AC}{\ifmmode\text{\texteuro}\else\texteuro\fi}
% Énoncés : Éducation nationale. Transcription et corrigés : APMEP, Denis Vergès et contributeurs. Découpage et classement : CoopMaths, Sébastien Lozano, Éric Elter et contributeurs. Adaptation de la présentation : Maths Explorer.
% CC BY-SA : https://coopmaths.fr/www/about/
% https://coopmaths.fr/alea/static/e3c/2020/tex/e3c_2020_00_sujet42_4_cor.tex
\begin{document}
\noindent\textit{APMEP / CoopMaths. CC BY-SA. Présentation adaptée par Maths Explorer.}\par\medskip
\subsection*{1.}

La tension est donnée par :
\begin{align*}
u(0) &= \sqrt{3} \sin\left(100\pi \times 0 + \dfrac{\pi}{3}\right) \\
&= \sqrt{3} \sin\left(\dfrac{\pi}{3}\right) \\
&= \sqrt{3} \times \dfrac{\sqrt{3}}{2} \\
&= \dfrac{3}{2} > \dfrac{\sqrt{3}}{2},
\end{align*}
la diode est donc passante à \( t = 0 \).

\subsection*{2.}

\begin{align*}
u\left(\dfrac{1}{100}\right) &= \sqrt{3} \sin\left(100\pi \times \dfrac{1}{100} + \dfrac{\pi}{3}\right) \\
&= \sqrt{3} \sin\left(\pi + \dfrac{\pi}{3}\right) \\
&= \sqrt{3} \times \left(-\sin \left(\dfrac{\pi}{3}\right)\right) \\
&= \sqrt{3} \times \left(- \dfrac{\sqrt{3}}{2}\right) \\
&= - \dfrac{3}{2} < \dfrac{\sqrt{3}}{2}.
\end{align*}
à \( t = \dfrac{1}{100} \), la diode n'est donc pas passante.

\subsection*{3.}

Puisque \( u\left(t + \dfrac{2}{100}\right) = u(t) \) pour tout \( t \geqslant 0 \), la fonction \( u \) est périodique de période \( T = \dfrac{2}{100} = 0{,}02 \) s.

\subsection*{4.}

\paragraph{a.} En observant la courbe, la diode semble devenir non passante autour de \( t = 0{,}005 \).

\begin{center}
\psset{xunit=500cm,yunit=1cm,algebraic=true,comma=true}
\begin{pspicture}(-0.0001,-2)(0.02,2)
\multido{\n=0.000+0.001}{21}{\psline[linewidth=0.2pt](\n,-2)(\n,2)}
\multido{\n=-2.0+0.5}{9}{\psline[linewidth=0.2pt](0,\n)(0.02,\n)}
\psaxes[linewidth=1.25pt,Dx=0.005,labelFontSize=\scriptstyle]{->}(0,0)(0,-2)(0.02,2)
\psplot[plotpoints=2000,linewidth=1.25pt,linecolor=blue]{0}{0.02}{1.73205*sin(314.159*x+1.0472)}


\psline[linewidth=1.25pt,linecolor=red]{->}(0.005,0.833)(0,0.833) 
\psline[linewidth=1.25pt,linecolor=red]{->}(0.005,0)(0.005,0.833) 
\uput[l](0,0.74){\red $\dfrac{\sqrt{3}}{2} \approx 0,83$} 

\end{pspicture}
\end{center}

\paragraph{b.}

\begin{align*}
u(0{,}005) &= \sqrt{3} \sin\left(100 \pi \times 0{,}005 + \dfrac{\pi}{3}\right) \\
&= \sqrt{3} \sin\left(\dfrac{\pi}{2} + \dfrac{\pi}{3}\right) \\
&= \sqrt{3} \sin\left(\dfrac{5\pi}{6}\right) \\
&= \sqrt{3} \times \frac{1}{2} \\
&= \frac{\sqrt{3}}{2}.
\end{align*}
La tension atteint exactement \( \dfrac{\sqrt{3}}{2} \) à \( t = 0{,}005 \), donc la diode devient non passante à cet instant.

\end{document}
